\documentclass[a4paper]{article}
\usepackage[utf8]{inputenc}
\usepackage{mathtools}
\usepackage[left=2cm,right=2cm,top=2cm,bottom=2cm]{geometry}
\usepackage{graphicx}
\usepackage{amsmath}
\usepackage{amssymb}
\usepackage{blindtext}
\usepackage{caption}
\usepackage{subcaption}
\usepackage{caption}
\usepackage{tabto}
\usepackage[export]{adjustbox}
\usepackage{float}
\usepackage{array}
\usepackage{booktabs}
\usepackage{svg}
\usepackage{comment}
\usepackage{siunitx}
\usepackage{braket}
\usepackage{hyperref}
\usepackage{physics}

\title{Important Calculations}
\author{Leart Zuka}
\date{\today}

\begin{document}


\maketitle
\tableofcontents
\newpage


\section{Phase Gate reflection differential equation stuff}
\subsection{Analytical derivation for case of atom in $\ket{0}$ and photon is in $\ket{h}$}

Our hamiltonian is as follows:
\begin{align*}
    H = \hbar g(\ket{e}\bra{1}a_h + \ket{1}\bra{e}a_h^{\dagger})
\end{align*}

We drive the $a_h$ mode of our cavity through $a_h^{in}(t)$, thus giving us the following equation of motion for our $a_h$ cavity mode: 

\begin{align*}
    \dot{a_h} &=-i[a_h,H]-(i\Delta + \frac{\kappa}{2})a_h - \sqrt{\kappa}a_h^{in}(t) \\
\end{align*}
Because $a_h$ commutes with $H$, the commutator is 0 and we get
\begin{align*}
    \rightarrow \dot{a_h}&=-(i\Delta + \frac{\kappa}{2})a_h - \sqrt{\kappa}a_h^{in}(t)
\end{align*}

We now drop the subscript $h$ as we'll only be looking at that case.

Multiplying with $e^{i\omega t}$, and integrating we get:
\begin{align*}
    \int{dt e^{i\omega t}\dot{a}(t)} =  \underbrace{[e^{i\omega t}a(t)]_{-\infty}^{\infty}}_{=0} - \underbrace{i\omega\int{dt e^{i\omega t}a(t)}}_{=i\omega a(\omega)}
\end{align*}
where in the last expression we used that the integral performs a Fourier transformation.

\begin{align*}
    \Rightarrow -i\omega a(\omega) &= -(i\Delta + \frac{\kappa}{2})a(\omega) - \sqrt{\kappa}a^{in}(\omega) \\
    \sqrt{\kappa}a^{in}(\omega) &= -(i\Delta + \frac{\kappa}{2} - i\omega)a(\omega) \\
    \rightarrow a(\omega) &= -\frac{\sqrt{\kappa}}{i\Delta+\frac{\kappa}{2}-i\omega}a^{in}(\omega)
\end{align*}

Now from our regular input output relation
\begin{align*}
    a^{out} (\omega)= a^{in}(\omega) + \kappa a(\omega) 
\end{align*}

We thus get:

\begin{align*}
    a^{out} (\omega)&=a^{in}(\omega) -\frac{\sqrt{\kappa}}{i\Delta+\frac{\kappa}{2}-i\omega}a^{in}(\omega) \\
    a^{out}(\omega) &=(1-\frac{\kappa}{i\Delta+\frac{\kappa}{2}-i\omega})a^{in}(\omega)\\
    a^{out}(\omega) &=(1-\frac{i\Delta + \frac{\kappa}{2}-i\omega-\kappa}{i\Delta+\frac{\kappa}{2}-i\omega})a^{in}(\omega)\\
    a^{out}(\omega) &=(\frac{i\Delta-\frac{\kappa}{2}-i\omega}{i\Delta+\frac{\kappa}{2}-i\omega})a^{in}(\omega)\\
    &=\underbrace{\frac{i(\Delta-\omega)-\frac{\kappa}{2}}{i(\Delta-\omega)+\frac{\kappa}{2}}}_{h(\omega)}a^{in}(\omega)
\end{align*}

We can then use the \href{https://en.wikipedia.org/wiki/Convolution_theorem}{convolution theorem} in order to solve for $a^{out}(t)$

\begin{align*}
    a^{out}(t) = \int_{-\infty}^{\infty}d\tau \underbrace{h(t-\tau)}_{h(t) =\frac{1}{2\pi}\int_{-\infty}^{\infty}d\omega e^{-i\omega t}h(\omega)}a^{in}(\tau)
\end{align*}

Now using that $h(\omega) = 1-\frac{\kappa}{i(\Delta-\omega)+\frac{\kappa}{2}}$ we get for $h(t)$ (via partial integration):

\begin{align*}
    h(t) &= \delta(t) - \kappa e^{-i\delta t}e^{-\frac{\kappa}{2}t}\theta(t) \\
    \Rightarrow a^{out}(t) &= a^{in}(t) - \kappa\int_{-\infty}^{\infty}d\tau e^{-i\delta(t-\tau)}e^{-\frac{\kappa}{2}(t-\tau)}a^{in}(t)
\end{align*}

For slowly varying pulses we can now approximate that $a^{in}(\tau)\approx a^{in}(t)$ thus giving us:

\begin{align*}
    \Rightarrow a^{out}(t) &\approx a^{in}(t) - \frac{\kappa}{i\Delta+\frac{\kappa}{2}}a^{in}(t)\\
    &=\frac{i\Delta-\frac{\kappa}{2}}{i\Delta+\frac{\kappa}{2}}a^{in}(t)
\end{align*}

Which for no detuning (so $\Delta=0$) gives us:
\begin{equation}
    a^{out}(t) =-a^{in}(t)
\end{equation}

\section{Calculation of Maxwell Bloch Equations \\(no decay to lossy channel)}

\subsection{Lindblad Formalism}

We start from the master equation:
\begin{equation}
\dot{\rho} = -\frac{i}{\hbar}[H,\rho] + \sum_j \gamma_j \left( L_j \rho L_j^\dagger - \tfrac{1}{2}\{ L_j^\dagger L_j, \rho \} \right),
\end{equation}
with
\begin{align}
\gamma_1 &= \gamma, \quad L_1 = \sigma, \\
\gamma_2 &= \kappa, \quad L_2 = a.
\end{align}

Now, for an operator \(O\),
\begin{align}
\frac{d}{dt} \langle O \rangle &= \Tr \left( O \dot{\rho} \right) \\
&= \Tr \left( O \left[ -\frac{i}{\hbar}[H,\rho] + \sum_j \gamma_j \left( L_j \rho L_j^\dagger - \tfrac{1}{2}\{ L_j^\dagger L_j, \rho \} \right) \right] \right).
\end{align}

The Hamiltonian contribution gives
\begin{equation}
- \Tr \left( O \frac{i}{\hbar}[H,\rho] \right)
= -\frac{i}{\hbar}\Tr \left( O H \rho - O \rho H \right)
= -\frac{i}{\hbar}\Tr \left( [O,H]\rho \right)
= -\frac{i}{\hbar} \langle [O,H] \rangle.
\end{equation}

For the dissipative terms we obtain
\begin{align}
\sum_j \gamma_j \Tr\Big( O L_j \rho L_j^\dagger 
- \tfrac{1}{2} O L_j^\dagger L_j \rho 
- \tfrac{1}{2} O \rho L_j^\dagger L_j \Big).
\end{align}

Using cyclicity of the trace:
\begin{align}
&= \sum_j \gamma_j \left( \langle L_j^\dagger O L_j \rangle 
- \tfrac{1}{2} \langle \{ L_j^\dagger L_j, O \} \rangle \right).
\end{align}

Therefore the general equation of motion is
\begin{equation}
\frac{d}{dt}\langle O \rangle 
= -\frac{i}{\hbar}\langle [O,H] \rangle 
+ \sum_j \gamma_j \left( \langle L_j^\dagger O L_j \rangle 
- \tfrac{1}{2}\langle \{L_j^\dagger L_j, O \} \rangle \right).
\end{equation}
\newpage
\subsection{Components of equation of motion}
\subsubsection{Jaynes Cummings Hamiltonian}
Assuming ground state has zero energy our Jaynes-Cummings Hamiltonian is equal to:
\begin{equation}
    H = \omega_ca^\dagger a + \omega_a\sigma^\dagger\sigma + g(\sigma + \sigma^\dagger)(a +  a^\dagger) + H_{drive}
\end{equation}
where $\omega_c$ is the atomic lowering operator, and $a$ is the photonic destruction operator.\\
$H_{drive}$ in this case is the external light field which is driving our system:
\begin{equation}
    H_{drive}^{I} = i(E(t)a^\dagger+E^*(t)a)
\end{equation}

\subsubsection{Cavity mode a}

From the Hamiltonian coupling term:
\begin{equation}
[ a, H ] = \hbar g (\sigma a^\dagger - \sigma^\dagger a).
\end{equation}
Thus

\begin{align}
\dot{a} &= -\frac{i}{\hbar}[a,H] - \frac{\kappa}{2}a \\
&= -ig\sigma - \frac{\kappa}{2}a.
\end{align}

\end{document}

